\originalchapter{转动惯量与静力平衡}\label{ch:rotation}
\originalsection{刚体定轴运动}
刚体绕固定轴转过同一个角 $\theta$, 但不同点到轴距离 $r_\perp$ 不同, 其速率为 $r_\perp|\omega|$. 把质点动能相加,
\begin{equation}
K=\tfrac12 I\omega^2,\qquad I=\sum_i m_ir_{\perp i}^2=\int r_\perp^2\dd m.
\end{equation}
$I$ 是关于指定轴的转动惯量, 单位 kg m$^2$. 同一物体换轴后 $I$ 一般改变. 质量远离轴会更难改变角速度, 平方权重反映这种分布效应.

均匀细杆长 $L$ 关于中点垂直杆轴, $I=(M/L)\int_{-L/2}^{L/2}x^2\dd x=ML^2/12$. 细圆环各点距轴均为 $R$, 所以 $I=MR^2$. 均匀实心圆盘用同心圆环积分, $\dd m=(2M/R^2)r\dd r$, 得 $I=MR^2/2$. 均匀球用垂直轴薄圆盘积分: 截面半径平方 $R^2-z^2$, $\dd I=\frac12\rho\pi(R^2-z^2)^2\dd z$, 从 $-R$ 到 $R$ 积分并代 $M=4\pi\rho R^3/3$, 得 $I=2MR^2/5$. 空心薄球壳相应为 $2MR^2/3$, 须区别实心球.

\begin{theorem}[平行轴定理]
距质心轴垂直距离 $d$ 的平行轴具有惯量 $I=I_C+Md^2$.
\end{theorem}
\begin{proof}
在质心平面坐标中令偏移沿 $x$ 方向. 展开 $\int[(x-d)^2+y^2]\dd m=I_C-2d\int x\dd m+Md^2$. 中间积分因质心坐标为零而消失. 不能用两条任意非质心轴直接套这个加法公式.
\end{proof}
平面薄片对穿过同一点的 $x,y,z$ 轴还满足 $I_z=I_x+I_y$, 因为 $r_{\perp z}^2=x^2+y^2$. 这是垂直轴定理, 限于平面薄片.

\originalsection{转动方程与功}
固定轴上力矩分量为 $\tau_z=(\vct r\times\vct F)\cdot\uv z$. 对恒定惯量的定轴转动, $\tau_z=I\alpha$. 可对各质点取切向运动方程并乘半径相加, 内部理想刚性约束的总内力矩抵消, 得到该式. 力矩做功 $\dd W=\tau_z\dd\theta$, 功率 $P=\tau_z\omega$. 若轴方向改变或 $I$ 随时间变化, 须用完整角动量方程.

有惯量 $I$ 的滑轮连接两块, 两侧张力一般不同. 重块 $m_2$ 向下加速度 $a$, 无滑条件 $\alpha=a/R$, 写
\begin{equation}
m_2g-T_2=m_2a,\quad T_1-m_1g=m_1a,\quad
(T_2-T_1)R=I a/R.
\end{equation}
消去张力得到 $a=(m_2-m_1)g/(m_1+m_2+I/R^2)$. 滑轮惯量相当于额外的惯性项, 不能把轮的全部质量直接加到平动质量中.

\originalsection{静力平衡与构件}
刚体静止需 $\sum\vct F=0$, 且 $\sum\vct\tau_O=0$. 合力为零时总力矩对不同参考点相同, 因为换点产生的差为两点位移叉乘总力. 选铰链为参考点可消去铰链力矩; 这并不意味着铰链力为零.

无质量直杆仅受两端铰链力时, 平衡要求二力等大、反向、共线, 称为二力杆. 拉伸杆在节点处向杆内部拉节点, 压缩杆则推节点远离杆. 桁架逐节点用水平和竖直平衡, 先从未知力少的端点入手. 若节点外有自重分布或中点负载, 该杆不再是二力杆.

\begin{example}
均匀杆质量 $M$, 长 $L$, 左端铰支, 右端用与杆夹角 $\beta$ 的绳拉住使杆水平. 求绳张力及铰支力.
\end{example}
\begin{solution}
关于左端力矩平衡: $TL\sin\beta=MgL/2$, 故 $T=Mg/(2\sin\beta)$. 竖直平衡给 $H_y=Mg-T\sin\beta=Mg/2$, 水平反力大小为 $T\cos\beta$. 绳几乎平行杆时张力会很大; 忽略角度不能估计安全载荷.
\end{solution}
\originalsection{重心与质心}
均匀重力场中总重力为 $M\vct g$, 关于任意点力矩为 $\int\vct r\times\vct g\dd m=\vct R\times M\vct g$. 所以可把总重量等效作用在质心, 此时重心与质心重合. 非均匀引力场中 $\vct g$ 随位置改变, 不能从积分中提出, 重心一般不再等于质心, 甚至不一定存在一个点可表示全部姿态下的重力矩. 本册杆与桁架静力题使用尺寸远小于地球半径的均匀场近似.
